\documentclass[11pt]{article}
\usepackage[utf8]{inputenc}
\usepackage[T1]{fontenc}
\usepackage{lmodern}
\usepackage[margin=0.85in]{geometry}
\usepackage{graphicx}
\usepackage[inkscape=false,inkscapelatex=false,inkscapepath=svgsubdir]{svg}
\usepackage{tikz}
\usetikzlibrary{positioning}
\usepackage{caption}
\usepackage{booktabs}
\usepackage{tabularx}
\usepackage{array}
\usepackage[table]{xcolor}
\usepackage{enumitem}
\usepackage{hyperref}
\usepackage{microtype}
\usepackage{parskip}

\hypersetup{colorlinks=true, linkcolor=black, urlcolor=blue, pdftitle={Washington's Crossing: Example of Play}}
\setcounter{tocdepth}{1}
\newcommand{\playimage}[2]{%
  \begin{figure}[htbp]
    \centering
    \includegraphics[width=0.9\linewidth,height=0.82\textheight,keepaspectratio]{#1}
    \caption{#2}
  \end{figure}%
}
\newcommand{\countergraphic}[2]{%
  \includesvg[width=#2]{../../Counters/Washingtons_Crossing/#1}%
}
\newcommand{\counterimage}[3]{%
  \begin{figure}[htbp]
    \centering
    \countergraphic{#1}{#2}
    \caption{#3}
  \end{figure}%
}
\newcommand{\smallplayimage}[2]{%
  \begin{figure}[htbp]
    \centering
    \includegraphics[height=0.45in,keepaspectratio]{#1}
    \caption{#2}
  \end{figure}%
}
\newcolumntype{Y}{>{\centering\arraybackslash}X}
\newenvironment{playchart}[1][0.78\linewidth]{%
  \begin{figure}[htbp]
    \centering
    \begin{minipage}{#1}
      \centering\small
      \renewcommand{\arraystretch}{1.16}
      \setlength{\tabcolsep}{6pt}
}{%
    \end{minipage}
  \end{figure}%
}
\newcommand{\charttitle}[1]{\textbf{\large #1}\par\smallskip}
\newcommand{\chartnote}[1]{\par\smallskip\begin{minipage}{\linewidth}\footnotesize\raggedright #1\end{minipage}\par}
\newcommand{\chartcaption}[1]{\par\medskip\caption{#1}}
\newcommand{\leadercounterdiagram}{%
  \begin{figure}[htbp]
    \centering
    \begin{tikzpicture}[font=\footnotesize, >=latex, counterarrow/.style={->,draw=gray!75,line width=0.45pt}]
      \node[inner sep=0] (commander) at (0,0) {\countergraphic{american-leader-cadwalader}{2.25cm}};
      \node[font=\bfseries, above=1.05cm of commander] {Commanding Leader};
      \node[anchor=east] (cost) at (-2.7,1.45) {Activation Cost};
      \draw[counterarrow] (cost.east) -- (commander.north west);
      \node[anchor=west] (span) at (2.2,1.45) {Command Span};
      \draw[counterarrow] (span.west) -- (commander.north east);
      \node[anchor=east] (rating) at (-2.7,-1.45) {Leader Rating};
      \draw[counterarrow] (rating.east) -- (commander.south west);
      \node[anchor=west] (radius) at (2.2,-1.45) {Command Radius};
      \draw[counterarrow] (radius.west) -- (commander.south east);
      \node[align=left,anchor=west] (militia) at (1.9,-0.22) {White box identifies\\Militia Leader};
      \draw[counterarrow] (militia.west) -- ([yshift=-0.22cm]commander.east);

      \node[inner sep=0] (subordinate) at (10.0,0) {\countergraphic{british-leader-rall}{2.25cm}};
      \node[font=\bfseries, above=1.05cm of subordinate] {Subordinate Leader};
      \node[anchor=east] (subcost) at (7.4,1.45) {Activation Cost};
      \draw[counterarrow] (subcost.east) -- (subordinate.north west);
      \node[anchor=east] (subrating) at (7.4,-1.45) {Leader Rating};
      \draw[counterarrow] (subrating.east) -- (subordinate.south west);
    \end{tikzpicture}
    \caption{Leader counters and their values.}
  \end{figure}%
}

\title{Washington's Crossing\\\large Campaign Game Example of Play}
\author{Roger Miller (original example)\\\small Collated from the locally archived Revolution Games web pages}
\date{2012 original material}

\begin{document}
\maketitle
\tableofcontents
\clearpage

\section{Introduction}

My name is Roger Miller. I am the game designer of \textit{Washington's Crossing}, the first game of our Campaigns of the American Revolution series. Since we are a new company and this is our first title, I felt it was necessary to introduce the game concepts while actually playing the game. This should hopefully give gamers the information they need to see if our game is for them.

\textit{Washington's Crossing} has two scenarios: the quick-start campaign, for those who wish to skip the initial Delaware River crossings by the Americans, and the Campaign Game, which I introduce here. The Campaign Game begins with General Washington preparing to cross the Delaware River to assault the Hessian garrison occupying Trenton, New Jersey. Before beginning, here is a brief overview of some game concepts.

\subsection{Leaders}

The primary maneuver pieces in the game are leaders. There are two leader displays in the game, one for the British player and one for the American player. Troop strengths and fatigue are tracked on the leader display. The number in parentheses next to Washington's name is the maximum number of troops he may have assigned to him: 140.

\counterimage{washington-leader-display}{0.9\linewidth}{Leader display.}

There are two types of leaders: those who can command other leaders (Commanding Leaders) and those who cannot (Subordinate Leaders).

\leadercounterdiagram

Leaders may move or attack only by spending activation points equal to their activation-point cost. Activated Commanding Leaders may activate other commanding and subordinate leaders. In this way, multiple leaders are activated by one overall commander. The following is an example of Washington's maximum command.

\playimage{Introduction_files/AmericanCommandExample_w7rE.JPG}{Example of Washington's maximum command.}

As leaders move and attack, their troops become fatigued. Fatigue is tracked on the leader display and ranges from 0 (fully rested) to 6 (extremely exhausted). Fatigue has negative effects on movement and combat.

\subsection{Movement}

Leaders do not have a fixed movement allowance. They roll a ten-sided die on the Movement Matrix chart to determine the number of movement points they have. Movement points can range from 0 to 10. The number across the top is the leader's rating; a leader with a rating of 5 uses the middle column. Several modifiers affect the die roll. Each time a leader moves, his troops gain 1 fatigue.

\begin{playchart}[0.72\linewidth]
  \charttitle{Movement Matrix (11.0)}
  \begin{tabular}{|>{\raggedright\arraybackslash}p{0.27\linewidth}|*{3}{>{\centering\arraybackslash}p{0.16\linewidth}|}}
    \hline
    \rowcolor{black!10}
    \textbf{\shortstack[l]{Modified\\Die Roll *}} & \textbf{2--4} & \textbf{5--6} & \textbf{7--8} \\
    \hline
    $-6$ or less & 0 MP & 0 MP & 0 MP \\ \hline
    $-5$ to $-3$ & 3 MP & 4 MP & 5 MP \\ \hline
    $-2$ to 0 & 4 MP & 5 MP & 6 MP \\ \hline
    1 to 3 & 5 MP & 6 MP & 7 MP \\ \hline
    4 to 6 & 6 MP & 7 MP & 8 MP \\ \hline
    7 to 9 & 7 MP & 8 MP & 9 MP \\ \hline
    10 or more & 8 MP & 9 MP & 10 MP \\ \hline
  \end{tabular}
  \par\vspace{1.35em}
  \textit{* Movement Matrix Die Roll Modifiers}\par\smallskip
  \begin{tabular}{|c|>{\centering\arraybackslash}p{0.34\linewidth}|c|>{\centering\arraybackslash}p{0.27\linewidth}|}
    \hline
    +4 & Night Move (American) & $-1$ & Disordered \\ \hline
    +2 & British Vanguard & $-1$ & Fatigue 1 \\ \hline
    $-2$ & All other British & $-2$ & Fatigue 2 \\ \hline
    $-1$ & Snow & $-3$ & Fatigue 3 \\ \hline
    $-2$ & Rain & $-4$ & Fatigue 4 \\ \hline
    +1 & Freeze & \multicolumn{2}{c|}{} \\ \hline
  \end{tabular}
  \chartcaption{Movement Matrix (11.0).}
\end{playchart}

\clearpage
\section{Game Setup}

The Campaign Game setup allows some variation in the American deployment. The British deployment is fixed. This gives the American player flexibility in planning the initial Delaware River crossings and strategy for attacking the starting British positions.

\counterimage{durham-boats}{0.28\linewidth}{Durham Boats setup marker (placed on a hexside next to hex 1330 or 1631).}

\subsection{American setup}

Right away, the American player must decide where to place his Durham Boats. He has two choices, marked by the red X: the hexside next to hex 1631, hoping to make the main crossing closer to Trenton; or the historical 1330, which is better for an initial attack on Princeton or Kingston. In this game, the American player places the Durham Boats in hex 1631 to support the crossing at Howell's and Yardley's Ferry.

\playimage{WC_GameSetup_files/DurhamBoatLocations_hY4C.JPG}{Possible Durham Boat locations.}

\playimage{WC_GameSetup_files/WashingtonSetup_hY4C.JPG}{Washington's initial setup.}

Washington, Sargent, Glover, Fermoy, St. Clair, and Lord Stirling all start within one hex of Newtown. All leaders except Lord Stirling are placed in 1434, planning to move to 1631 and cross the river at Howell's and Yardley's Ferry with the assistance of the Durham Boats (red arrow). Lord Stirling is placed in 1334, planning to move to 1330 and cross the river at Johnson's and McConkey's Ferry (red arrow).

\playimage{WC_GameSetup_files/Deployment_1_hY4C.JPG}{Initial American deployment.}

Setting up to cross the Delaware River in several locations is a good plan to reduce the risk of poor river-crossing die rolls. Sullivan and Greene, who start anywhere in Pennsylvania, are placed in 1332 and 1633 respectively. This is within Washington's command range and within range of all the American leaders around Newtown, including Stephen and Mercer, located next to the river in 1330 and 1631. This will allow Washington to activate the entire group, saving activation points and allowing all leaders to move using Washington's Leader Rating.

\subsection{Initial strategy}

The American player has three turns---two night turns and one day turn---before the British player moves. This gives the American player the potential to cover quite a bit of distance to set up attacks. Potential targets are Rall at Trenton, Leslie at Princeton, Light Infantry at Kingston, Hessian Detachment 2 at Bordentown, von Donop at Mt. Holly, and Hessian Detachment 1 in the marsh near Bordentown.

The American player has the strength to attack multiple locations on the first day turn and should attack Trenton and one or two other locations if possible. The main problem is the weather and its potential to negatively affect crossing the Delaware River. The plan must remain flexible. It is great to plan to cross at Burlington with Cadwalader's force and march on Bordentown to attack von Donop, but if the weather is bad and the crossing die roll is poor, the plan must be adjusted. Historically, Cadwalader got only 700 men across and no artillery on the first day; he called off his crossing and his part in the Trenton attack. He rested his men and crossed a day later, leading to the Princeton campaign when Washington decided to follow up Cadwalader's successful crossing.

Finally, if you decide to attack Princeton or Kingston, you will cross the orange line across the map, triggering accelerated British reinforcements. You had better have a plan to follow up the initial attacks with an aggressive move north to take the main British towns or kill a lot of British troops. This strategy is not recommended for new players, but for an experienced player it is viable and a lot of fun.

In this game, the American player plans to attack Trenton with a surrounding attack from Washington's force and to assault Bordentown with Cadwalader's force. The British player has no initial strategy; he must wait to see where the Americans move and how they attack before making a plan.

\playimage{WC_GameSetup_files/Deployment2_hY4C.JPG}{American plan and British starting positions.}

\clearpage
\section{Turn 1: December 25, 1776, 5 p.m. (Night Turn)}

\begin{playchart}[0.76\linewidth]
  \charttitle{Sequence of Play}
  \raggedright
  \begin{enumerate}[label=\arabic*., leftmargin=1.9em, itemsep=0.2em, topsep=0.1em]
    \item Weather Phase
    \item American Raid Phase
    \item Reinforcement Phase
    \item American Command Phase
      \begin{enumerate}[label=\alph*., leftmargin=1.8em, itemsep=0pt, topsep=0pt]
        \item American player activates leaders with Activation Points.
      \end{enumerate}
    \item American Action Phase
      \begin{enumerate}[label=\alph*., leftmargin=1.8em, itemsep=0pt, topsep=0pt]
        \item Move and perform Normal/Hasty attacks with leaders activated by activation points (Defender may react; see 11.3 Reaction Movement).
        \item Move and perform Normal/Hasty attacks with leaders activated by initiative roll (one leader per day turn) (Defender may react; see 11.3 Reaction Movement).
        \item Conduct all Prepared attacks (Defender may react; see 11.3 Reaction Movement).
      \end{enumerate}
    \item American Recovery Phase \textit{(Night Turns Only)}
    \item British Command Phase
      \begin{enumerate}[label=\alph*., leftmargin=1.8em, itemsep=0pt, topsep=0pt]
        \item British player activates leaders with Activation Points.
      \end{enumerate}
    \item British Action Phase
      \begin{enumerate}[label=\alph*., leftmargin=1.8em, itemsep=0pt, topsep=0pt]
        \item Move and perform Normal/Hasty attacks with leaders activated by activation points (Defender may react; see 11.3 Reaction Movement).
        \item Move and perform Normal/Hasty attacks with leaders activated by initiative roll (one leader per day turn) (Defender may react; see 11.3 Reaction Movement).
        \item Conduct all Prepared attacks (Defender may react; see 11.3 Reaction Movement).
      \end{enumerate}
    \item British Recovery Phase \textit{(Night Turns Only)}
    \item End of Turn
  \end{enumerate}
  \chartcaption{Sequence of play.}
\end{playchart}

\subsection{Weather Phase}

The first thing that happens during each turn is weather determination. Weather determination in \textit{Washington's Crossing} is a two-step process. Players first check whether the previous turn's weather continues for this turn: roll one die and consult the Continuing Weather table.

\begin{playchart}[0.72\linewidth]
  \charttitle{Continuing Weather (7.0)}
  \begin{tabular}{|>{\raggedright\arraybackslash}p{0.19\linewidth}|>{\centering\arraybackslash}p{0.69\linewidth}|}
    \hline
    \rowcolor{black!10}\textbf{Die Roll} & \textbf{Result} \\
    \hline
    0--2 & Continues \\ \hline
    3 & Continues if Freeze, otherwise New Weather \\ \hline
    4--5 & Continues if Rain, otherwise New Weather \\ \hline
    6--8 & Continues if Snow, otherwise New Weather \\ \hline
    9 & New Weather (see table below) \\ \hline
  \end{tabular}
  \chartcaption{Continuing Weather table (7.0).}
\end{playchart}

If the weather does not continue, players consult the New Weather Table to check what the current turn's weather will be.

\begin{playchart}[0.72\linewidth]
  \charttitle{New Weather (7.0)}
  \chartnote{If night, subtract one from the die roll.}
  \begin{tabular}{|>{\raggedright\arraybackslash}p{0.19\linewidth}|>{\centering\arraybackslash}p{0.69\linewidth}|}
    \hline
    \rowcolor{black!10}\textbf{Die Roll} & \textbf{Result} \\
    \hline
    $-1$ & Freeze \\ \hline
    0 & Snow \\ \hline
    1--4 & Freeze \\ \hline
    5 & Rain \\ \hline
    6--9 & Fair \\ \hline
  \end{tabular}
  \chartcaption{New Weather table (7.0).}
\end{playchart}

For the first game turn, the previous turn's weather is Fair. The American player prays that this continues in the hope of an easy crossing of the Delaware River. A die roll of 4 indicates that the previous turn's weather would continue if it was Rain; since the previous turn was not Rain, move on to the New Weather table. A die roll of 2, modified down 1 because it is a night turn, yields Freeze.

\subsection{American Raid Phase}

The American player rolls a 2 on the Raid Table, which indicates that a raid is possible. He rolls on the Raid Leader Table to see which leader conducts the raid, and rolls a 7: Stephen's brigade. Stephen may choose any British leader within 15 hexes as his target, and chooses the Hessian leader Rall in Trenton, trying to soften him up before Washington's attack. A die roll of 0, however, causes Stephen to lose 20 men and does no damage to Rall other than 1 fatigue level, since all raids at night increase the target's fatigue by 1. Rall is now at fatigue level 3.

\begin{playchart}[0.68\linewidth]
  \charttitle{A. Raid Table (8.0)}
  \begin{tabular}{|>{\raggedright\arraybackslash}p{0.20\linewidth}|>{\centering\arraybackslash}p{0.70\linewidth}|}
    \hline
    \rowcolor{black!10}\textbf{Die Roll} & \textbf{Result} \\
    \hline
    0--2 & Raid --- consult Raid Leader Table \\ \hline
    3--6 & No effect \\ \hline
  \end{tabular}
  \par\medskip
  \charttitle{B. Raid Leader Table (8.0)}
  \begin{tabular}{|>{\raggedright\arraybackslash}p{0.20\linewidth}|>{\centering\arraybackslash}p{0.70\linewidth}|}
    \hline
    \rowcolor{black!10}\textbf{Die Roll} & \textbf{Leader} \\
    \hline
    0--1 & Griffin \\ \hline
    2--3 & Ewing \\ \hline
    4--5 & Dickinson \\ \hline
    6--7 & Stephen \\ \hline
    8--9 & McDougall \\ \hline
  \end{tabular}
  \chartcaption{Raid and Raid Leader tables (8.0).}
\end{playchart}
\begin{playchart}[0.75\linewidth]
  \charttitle{C. Raid Result Table (8.0)}
  \chartnote{On night turns, the defending leader receives one additional fatigue on all results.}
  \begin{tabular}{|>{\raggedright\arraybackslash}p{0.14\linewidth}|>{\centering\arraybackslash}p{0.76\linewidth}|}
    \hline
    \rowcolor{black!10}\textbf{Die Roll} & \textbf{Result} \\
    \hline
    0 & Attacker loses 20 troops \\ \hline
    1--2 & Attacker loses 10 troops \\ \hline
    3 & No effect \\ \hline
    4--5 & Defender loses 10 troops and takes 1 Fatigue \\ \hline
    6 & Defender loses 20 troops and takes 1 Fatigue \\ \hline
    7 & Defender loses 20 troops and the British player loses 1 Activation Point \\ \hline
    8 & Defender loses 30 troops and takes 1 Fatigue \\ \hline
    9 & Defender loses 30 troops and the British player loses 1 Activation Point \\ \hline
  \end{tabular}
  \chartcaption{Raid Result Table (8.0).}
\end{playchart}

\subsection{Reinforcement Phase}

There are no reinforcements on game turn 1, so move on to the next phase.

\subsection{American Command Phase}

The American player begins with 21 activation points and decides which leaders to activate. He activates Washington for 4 activation points. Washington has a printed command span of 6, and activates Greene and Sullivan; they in turn use their command spans to activate Stephen, Lord Stirling, St. Clair, Mercer, Sargent, Fermoy, and Glover.

\subsection{American Action Phase}

The American player will now move and attack with his leaders. Since this is a night turn, only movement is allowed. Night is the ideal time for the Americans to move: they receive a +4 modifier to their movement die rolls.

The American player starts by moving Stephen. To find out how far Stephen can move, he rolls one die and consults the Movement Matrix chart. Since Stephen was activated by Washington---even though activation first went to Sullivan and then Stephen---he uses the 7--8 column representing Washington's Leader Rating. A die roll of 9 is modified up to 10, the maximum on the chart, because Freeze is a +1 modifier and night is a +4 modifier for the Americans. Stephen receives 10 movement points and tries to cross the Delaware River into hex 1329. The attempt costs 4 movement points and requires a die roll on the Delaware River Crossing Table. A roll of 3, modified down 2 for Freeze and up 1 because it is before December 28, yields 2: no troops may cross. Stephen could move elsewhere, but ends his move and hopes for a better crossing die roll next turn. His fatigue increases by 1, to level 1.

The American player then decides to move Washington and all leaders next to him. Since they are all adjacent, he rolls once to determine the movement points for all leaders except Washington, who always moves 13. A die roll of 8 is modified by night and Freeze to 10, and all of Washington's brigade commanders have 10 movement points. Washington and Glover move to 1631, spending 4 movement points, and attempt to cross the Delaware River for another 4. The crossing die roll is 5, modified down 2 for Freeze; up 2 for the Durham Boats; up 1 for Washington; up 1 for Glover; and up 1 for being before December 28. This produces an 8 on the two-Ferry column and allows 2,500 troops to cross this hexside during this Action Phase. Washington and Glover pay 2 more movement points to enter 1630 and 1629. Washington and Glover attempted the crossing because they have positive modifiers.

Sargent then moves 10 movement points to 1629. He need not check the river crossing, because the player knows that 2,500 troops may cross this turn and Washington, Glover, and Sargent add up to 1,700 troops. Lord Stirling, Fermoy, and St. Clair move to 1631, planning to cross next turn. Mercer determines his movement points and gets 9. He crosses the river, using 720 more troops of the remaining 800 available to cross, and moves to 1728. All American leaders with troops that moved gain 1 fatigue, to level 1. Greene and Sullivan now move to hex 1630. They have no troops assigned and always move 13 movement points.

\playimage{WC_turn1_files/Turn1_Situation_JvYc.JPG}{Situation after Turn 1.}

\subsection{Recovery Phase}

Recovery occurs only during night turns. All leaders who did not move or entrench may regain 1 fatigue level. The only American leader who meets this condition is Detachment 4, which reduces its fatigue from 1 to 0 because it did not move.

\subsection{British Turn}

The British player is not yet allowed to move any units, so his only action is recovery from fatigue. Rall reduces his fatigue by 1, from 3 to 2.

\clearpage
\section{Turn 2: December 26, 1776, 12 a.m. (Night Turn)}

\subsection{Weather Phase}

The Continuing Weather die roll is a 9: roll for new weather. The New Weather roll is 6, modified down 1 for night, so the weather changes to Rain. This is a break for the Americans: Freeze and Snow are negative river-crossing modifiers, while Rain only slows movement somewhat.

\subsection{American Raid and Reinforcement Phases}

A die roll of 5 means no raid this turn. There are no reinforcements for either side.

\subsection{American Command Phase}

The American player receives 8 activation points, as in every 12 a.m. turn except December 30, when he receives none. He now has 25 activation points. He activates Washington for 4 activation points; Washington activates Mercer and Sullivan with his command span, and Sullivan uses his command span to activate Lord Stirling, St. Clair, and Fermoy. The American player also activates Stephen, Dickinson, and Ewing for 1 activation point each, and Cadwalader for 2 activation points; Cadwalader uses his command span to activate Hitchcock and Detachment 7. The player has spent 9 activation points and has 16 left.

\subsection{American Action Phase}

The American player starts by rolling for Stephen's movement. This turn Stephen uses his own leadership value of 4 because he was not activated by another leader. He receives 6 movement points and again tries to cross the Delaware River to hex 1329. This time he rolls a 6, and a maximum of 2,000 troops may cross the Ferry. Stephen moves to 1329, using 4 movement points for the Ferry crossing, 1 for hex 1329, and 1 for 1328.

\playimage{WC_Turn2_files/Stephen_Crossing_haxh.JPG}{Stephen's crossing.}

Dickinson rolls for movement, gets 6 movement points, and tries to cross the Delaware River to hex 1930. His modified 8 on the one-Ferry column allows 1,250 troops to cross. Dickinson has only 400 troops, so he enters 1930, using 4 movement points for the Ferry, 1 for the hex, and 1 for Rall's Zone of Control in Trenton.

Ewing rolls 8 movement points and attempts to cross to hex 2130. A modified 5 allows 500 troops to cross, but Ewing has 720 troops. The rules require the American player to cross as many troops as possible from the unit attempting the crossing. Ewing crosses with 500 and leaves Detachment Militia No. 5, with 220 troops, behind. Ewing ends in 2130, cutting off Rall's retreat and remaining next to the Ferry in the hope of getting his remaining troops across in later turns. The detachment has also finished moving, because when a detachment is created during movement, either the parent unit or the detachment must cease moving. Ewing continued, so the detachment is done moving.

\playimage{WC_Turn2_files/EwingMove_haxh.JPG}{Ewing's move and Detachment Militia No. 5.}

St. Clair, Lord Stirling, and Fermoy roll together because they are stacked and activated by the same leader, receiving 9 movement points. Fermoy tries to cross to hex 1630; a modified 6 allows 1,500 troops to cross, and Fermoy has 620. Fermoy continues to 1629. It is a good idea to check the river crossing with a single leader so that, if it goes poorly, other leaders have not spent the 4 movement points for a crossing attempt or fatigued themselves. Lord Stirling, with 570 troops, also moves to 1629. St. Clair, with 420 troops, cannot get all his troops across because only 310 capacity remains, so he chooses not to move. Mercer rolls 9 movement points and moves to 2129, cutting off another hex of Rall's potential retreat.

Cadwalader rolls with his group and gets 7 movement points. His crossing result allows 1,750 troops, permitting Cadwalader and Hitchcock to cross and move to 2537. Detachment 7 also succeeds and moves to 2537. Sullivan always has 13 movement points and moves to 1629 to join Washington. Greene remains in place to keep St. Clair in his command span. All American units that moved gain 1 fatigue. Washington and St. Clair were activated but did not move, so they gain no fatigue and may recover. Detachment 5 was created after Ewing spent 4 movement points on the crossing attempt, but gains fatigue as it has moved.

\subsection{American Recovery Phase}

All fatigued American leaders who did not move regain one fatigue level. Washington, St. Clair, Glover, and Sargent each recover 1, all to zero. Cadwalader, Hitchcock, Detachment 7, Dickinson, Detachment 5, and Ewing are at fatigue 1. Mercer, Stephen, Fermoy, and Lord Stirling are at fatigue 2.

\subsection{British Turn}

The British still cannot move in Turn 2. They receive 8 activation points in their Command Phase, as they do every 12 a.m. turn, for a total of 8 because they began the game with none. In the British Recovery Phase, Rall recovers one fatigue level, to level 1.

\playimage{WC_Turn2_files/Turn2_Situation_haxh.JPG}{Situation at the end of Turn 2.}

\playimage{WC_Turn2_files/Turn2_Cdwldr_haxh.JPG}{Cadwalader and von Donop.}

\clearpage
\section{Turn 3: December 26, 1776, 7 a.m. (Day Turn)}

\subsection{Weather, Raid, and Reinforcement Phases}

Weather changes to Fair. There is no raid and no reinforcements for either side.

\subsection{American Command Phase}

Cadwalader absorbs the 350 troops from Detachment 7, which is removed from the map for reuse. Cadwalader does not need to be activated to absorb the detachment. The American player activates Washington for 4 activation points, then activates Stephen, Dickinson, and Greene with his command span. Greene activates Glover, Sargent, Lord Stirling, Fermoy, and St. Clair with his command span. The player activates Cadwalader for 2 activation points, and Cadwalader activates Hitchcock with his command span. The American player has 10 activation points remaining.

\subsection{American Action Phase}

Cadwalader and Hitchcock move to 2831 and increase their fatigue to 2. Stephen is far from the attack on Trenton planned for this turn, so the player declares that Stephen will force march. Stephen rolls 6 movement points. Force marching adds 3 (only 2 for a British leader), so he can move 9 and moves to 1930. Stephen gains 1 fatigue for the move and 1 for force marching, and is now at fatigue 4.

\playimage{WC_Turn3_files/turn3_move_SpxH.JPG}{Turn 3 movement; Stephen's move is shown in white.}

Dickinson rolls 7 movement points. Starting in 1930 (movement in yellow), he enters 1929 for 3 movement points: 1 for the road, 1 to leave a Zone of Control in 1930, and 1 to enter a Zone of Control in 1929. Dickinson continues to 2028, spending 4 more, and reaches fatigue 2. St. Clair rolls 9 movement points and tries to cross the Ferry to 1630. A die roll of 0, modified by the Durham Boats and date to 3, allows no troops to cross. He has spent 4 movement points on the attempt, so moves downriver to 2030, spending 4 more; his fatigue is now 1.

Glover, Sargent, Lord Stirling, and Fermoy roll together and receive 6 movement points. Washington moves with them and has 13. They move from 1629 to 1729, 1828, and 1929, spending 3 for road movement and 1 to enter Rall's Zone of Control in Trenton. The American player attacks Trenton with Washington, paying 1 movement point for the terrain Rall occupies and 1 for a hasty attack. With only 6 movement points, this is the best attack available.

\subsection{Attack on Trenton}

Washington's stack has 2,890 troops attacking Rall's 1,250. First determine Surprise modifiers. The American player rolls 6 on the Surprise table, then applies -1 for the hasty attack and +7 for attacking Trenton in the game's first a.m. turn---``think Pearl Harbor surprise''---for a result of 12. The American receives +4 to the combat die roll and -10\% casualties from surprise.

\begin{playchart}[0.72\linewidth]
  \charttitle{Surprise Table (17.1)}
  \begin{tabular}{|>{\raggedright\arraybackslash}p{0.38\linewidth}|>{\centering\arraybackslash}p{0.25\linewidth}|>{\centering\arraybackslash}p{0.25\linewidth}|}
    \hline
    \rowcolor{black!10}\textbf{\shortstack[l]{Modified\\Die Roll *}} & \textbf{CRT Modifier} & \textbf{Attacker \%} \\
    \hline
    $-4$ or less & $-6$ & +15 \\ \hline
    $-3$ & $-5$ & +10 \\ \hline
    $-2$ & $-4$ & +10 \\ \hline
    $-1$ & $-3$ & +10 \\ \hline
    0--1 & $-2$ & +5 \\ \hline
    2 & $-1$ & +5 \\ \hline
    3 & $-1$ & --- \\ \hline
    4--5 & 0 & --- \\ \hline
    6--7 & +1 & --- \\ \hline
    8--9 & +2 & $-5$ \\ \hline
    10 & +3 & $-5$ \\ \hline
    11 & +3 & $-10$ \\ \hline
    12 & +4 & $-10$ \\ \hline
    13 & +5 & $-10$ \\ \hline
    14 or more & +6 & $-15$ \\ \hline
  \end{tabular}
  \chartcaption{Surprise Table (17.1).}
\end{playchart}

The American player rolls on the Combat Result Table using the 7--8 column because Washington has a Leader Rating of 8. He rolls 2, modified -1 for the hasty attack, +1 for Americans attacking, +4 for surprise, +2 for outnumbering Rall 2 to 1, and +4 because American leaders occupy all six hexes around Rall. The final result is 12: 20/60! The attacker result is on the left, so Washington would take 20\% of the defender's strength in casualties. Surprise decreases Washington's casualties by 10\%, so the Americans lose 10\% of Rall's strength: 130 troops. The player distributes all of these losses to Lord Stirling, reducing him from 570 to 440 troops.

Rall takes 60\% of 1,250, or 750 casualties, leaving 500. The exclamation mark means that Rall routs and must retreat four hexes. Completely surrounded, Rall chooses to retreat through Dickinson in 2028, then to 2027, 1927, and 1926. Retreating through Dickinson causes another 10\% loss of Rall's original strength, 130 more casualties, reducing him to 370. When surrounded, the first retreat hex is up to the retreating force; after that, retreat priorities dictate the path: do not retreat through enemy leaders with troops, avoid enemy Zones of Control, retreat toward the map edge (the top edge for the British), and stay on roads if possible. Rall gains 3 fatigue: 1 for defending in combat and taking at least 5\% loss, and 2 for routing. He is now at fatigue 4.

Washington must advance after combat into Trenton and may continue moving with troops that still have movement points. The other leaders with troops rolled 6 and have used all 6, so only Washington and Greene, who always have 13, can keep moving. The American player is content with their positions and halts in Trenton. Washington, Glover, Sargent, Lord Stirling, and Fermoy gain 2 fatigue: 1 for moving and 1 for attacking. Washington, Glover, and Sargent are at 2; Fermoy and Lord Stirling are at 4. Greene has no troops assigned and gains none. The American player earns 6 victory points for a major victory. Sullivan moves to 2130. Since he has no troops assigned and activated no other leaders, he may activate for free and move 13 without fatigue.

\playimage{WC_Turn3_files/Rall_attack_SpxH.JPG}{Attack on Rall at Trenton.}

The American player may now activate by die roll one leader with an activation cost of 2 or less, on day turns only. He tries Detachment 4, with leadership rating 3, so must roll 3 or less. A roll of 2 activates the detachment. It rolls 7 movement points, moves to 2346, and reaches fatigue 1.

\subsection{American Recovery Phase}

There is no recovery during day turns, so this phase is skipped. A leader such as Stephen, at fatigue 4, needs four night turns---or two complete nights---to recover to zero. He cannot move or attack in the day turns without increasing fatigue. This cycle of action and rest is central to the game and as important to manage well as maneuver or combat.

\subsection{British Command and Action Phases}

On this turn the British may activate only certain leaders; the rest are still thinking that nothing will happen until spring. Rall may move without activation because he is routed and within 10 movement points of enemy troops. von Donop, Stirling, Hessian Detachment 1, and Hessian Detachment 2 are available. The British have only 8 activation points and conservatively spend 2 to activate the two Hessian detachments.

When rolling movement for British leaders, several facts become apparent. The average British leader rolls on a worse movement column than the Americans, as British Leader Ratings are generally lower (with exceptions, such as the American Fermoy). All British units except Vanguard have a -2 movement-die modifier. British force marches gain only 2 movement points to the American 3. The British therefore move more slowly and must plan much further ahead or the Americans will run circles around them.

Hessian Detachment 2 rolls 5 and moves to 2927; Hessian Detachment 1 rolls 4 and moves to 2827. Both gain 1 fatigue. Rall gets 10 movement points while routed and moves to 1816, trying to move away from enemy troops and toward his supply base. He gains 1 fatigue and is now at 5. The British may try to activate one leader by die roll and attempt von Donop. A roll of 3, lower than von Donop's Leader Rating of 5, activates him; he uses his command span to activate Stirling. von Donop rolls group movement with Stirling, gets 5, and both move to 2641. Rall and Stirling both gain 1 fatigue. There is no British recovery because it is a day turn.

\playimage{WC_Turn3_files/Turn3_NorthFront_SpxH.JPG}{Situation at the end of Turn 3: North Front.}
\playimage{WC_Turn3_files/Turn3__CentralFront_SpxH.JPG}{Situation at the end of Turn 3: Central Front.}
\playimage{WC_Turn3_files/Turn3_sFront_SpxH.JPG}{Situation at the end of Turn 3: South Front.}

\clearpage
\section{Turn 4: December 26, 1776, 12 p.m. (Day Turn)}

\subsection{Weather, Raid, and Reinforcement Phases}

Weather is Fair. There is no American raid and no reinforcements for either side.

\subsection{American Turn}

The American player spends 1 point to activate Griffin and 3 to activate Greene. Greene uses his command span to activate St. Clair, Dickinson, and Detachment 5. The American has 6 activation points remaining.

St. Clair moves to 2029, Sullivan to 2029, Dickinson to 2329, Greene to 2329, Griffin to 2346, and Detachment 5 to 2130.

\subsection{British Turn}

The British player spends 1 point to activate Hessian Detachment 1, 1 to activate Hessian Detachment 2, and 2 to activate von Donop, who uses his command span to activate Stirling. Hessian Detachment 1 moves to 2923, Hessian Detachment 2 to 2821, and von Donop and Stirling to 2445. Leslie tries to activate by die roll and fails. Rall recovers from rout because he is more than 10 movement points from the enemy; he is now disordered.

\subsection{Commentary}

The American army is waiting for night to rest, moving only leaders who are lightly fatigued or were on the wrong side of the river and needed to cross while weather is good. One exception is the area where Griffin joined Detachment 4 and the detachment started digging entrenchments. These leaders feel threatened by von Donop's aggressive move.

The British are moving the Hessian detachments north to put distance between themselves and Washington's main force. von Donop is trying to attack one of the smaller American forces before it can be reinforced or dig in. Now that Griffin has met Detachment 4 and begun to dig, both sides must decide whether this is the time and place for a battle on the next morning.

\playimage{WC_turn4_files/North_Turn4_HlS6.JPG}{Situation after Turn 4: British northern positions.}
\playimage{WC_turn4_files/Central_Turn4_HlS6.JPG}{Situation after Turn 4: American central position.}
\playimage{WC_turn4_files/southern_turn4_HlS6.JPG}{Situation after Turn 4: Southern Front.}

\clearpage
\section{Turn 5: December 26, 1776, 7 p.m. (Night Turn)}

\subsection{Weather, Raid, and Reinforcement Phases}

Weather is Fair. There is no American raid and no reinforcements for either side.

\smallplayimage{WC_Turn5_files/FairWeatherSymbol_FQ5g.JPG}{Fair weather symbol.}

\subsection{American Turn}

Ewing absorbs Detachment 5. All troops are moved to Ewing, increasing his strength to 720, the maximum he may have assigned. Detachment 5 had fatigue 2, higher than Ewing's fatigue 1, so Ewing now has fatigue 2. Detachment 5 is removed from the map for reuse. The American player activates no leaders, so there is no American Action Phase.

During the American Recovery Phase, all American leaders recover 1 fatigue if they have not moved or entrenched.

\subsection{British Turn}

The British may activate a maximum of one leader each night turn. The British player activates Leslie for 1 activation point. Leslie's activation cost is 2 if he uses his command span to activate other leaders and 1 if he does not. The British now have 1 activation point remaining. Leslie moves to 1818.

During British Recovery, all British leaders recover 1 fatigue if they did not move or entrench. Rall also removes his disorder.

\subsection{Commentary}

If you wish to move a short distance during the night and were not previously fatigued, it is wise to move in the first night turn and rest in the second. Your troops will start the next day with no fatigue.

\clearpage
\section{Turn 6: December 27, 1776, 12 a.m. (Night Turn)}

\subsection{Weather, Raid, and Reinforcement Phases}

Weather is Fair. There is no American raid and no reinforcements for either side.

\smallplayimage{WC_Turn6_files/FairWeatherSymbol_Cvqs.JPG}{Fair weather symbol.}

\subsection{American Turn}

The American player gains 8 activation points, as in every 12 a.m. turn except December 30, when he gains none. He now has 14 activation points, but activates no leaders, so there is no Action Phase.

During American Recovery, all leaders who did not move or entrench recover 1 fatigue. All American leaders are at fatigue 0 except Stephen, Lord Stirling, and Fermoy, all at fatigue 2, and Dickinson at fatigue 1.

\subsection{British Turn}

The British player gains 8 activation points, as in every 12 a.m. turn, and now has 9. He activates no leaders, so there is no Action Phase. During British Recovery, all leaders who did not move or entrench recover 1 fatigue. All British leaders are at fatigue 0 except Rall, who is at fatigue 3.

\subsection{Commentary}

Night turns are generally very quick because most of the time both players are resting. However, the American player should always look to take advantage of superior night mobility: +4 to the Movement Matrix die roll. The British have very limited night moves: one leader with an activation cost of 2 or less. Rall is still at fatigue 3. A major battle that causes a leader to rout usually piles up so much fatigue that it takes several days of not moving and several nights of resting before the unit regains full effectiveness.

\vfill
\begin{center}
\small Original text and images: \copyright\ 2012 Revolution Games. This document is a local collation of the downloaded example-of-play pages, excluding website navigation and decorative elements.
\end{center}

\end{document}
